\chapter{Inferior Vena Cava Filter Placement}

\section*{Overview}
Inferior vena cava (IVC) filter placement is an interventional procedure in which a metal filter is deployed in the inferior vena cava to trap emboli traveling from the lower extremities. It is a mechanical means of preventing pulmonary embolism.

\section*{Alternative Names}
IVC filter, vena cava filter, Greenfield filter, retrievable vena cava filter

\section*{Principle}
The filter is positioned in the infrarenal IVC, where it spans the lumen and captures thrombus migrating from pelvic or lower extremity veins while allowing normal blood flow. Trapped clot is subsequently broken down by the body's endogenous fibrinolytic system, and the filter is removed once the risk of embolism resolves.

\section*{History}
The concept of caval interruption dates to the 1960s, and the first percutaneously placed IVC filter was introduced in the 1970s. Retrievable filters developed in the 1990s allow removal once protection is no longer needed.

\section*{Why It Is Done}
Placed in patients with proven venous thromboembolism and a contraindication to anticoagulation, failure or complications of anticoagulation, or recurrent pulmonary embolism despite adequate therapy. It is also considered for prophylaxis in select high-risk trauma patients or those undergoing major surgery with a high bleeding risk.

\section*{Is This a Primary or Secondary Test or Procedure}
A secondary procedure to prevent pulmonary embolism when pharmacologic anticoagulation is contraindicated or has failed, rather than a primary treatment of the thrombus itself.

\section*{How It Is Performed}
Venous access is obtained, usually through the femoral or internal jugular vein, and a venogram or ultrasound confirms IVC anatomy and the level of the renal veins. A sheath containing the collapsed filter is advanced to the infrarenal IVC, the filter is deployed, and a completion venogram confirms position and patency.

\section*{Preparation}
Informed consent, imaging to assess IVC size, diameter, and anatomy, and assessment of renal function before iodinated contrast venography. Anticoagulation is held or reversed as appropriate for the bleeding risk.

\section*{Contraindications and Precautions}
Contraindications include an IVC diameter exceeding the filter's capacity, severe thrombus in the IVC, and inability to access the venous system; caution is needed in pregnancy. Precautions include documenting the filter in registries, planning timely retrieval, and avoiding prolonged placement because of the risk of filter thrombosis and fracture.

\section*{Risks}
Filter thrombosis or IVC occlusion, filter migration or tilting, strut fracture, vein injury or bleeding at the access site, and complications of retrieval including perforation.

\section*{Aftercare and Recovery}
The patient resumes anticoagulation once safe and undergoes regular follow-up imaging to assess filter position. Retrieval is attempted when the indication for protection has resolved, ideally within weeks to months.

\section*{Outcome and Follow-up}
A completion venogram showing the filter centered in the infrarenal IVC below the renal veins with brisk flow indicates successful placement. Imaging at follow-up showing stable filter position, patent IVC, and no new pulmonary emboli reflects a good outcome.

\section*{Expected Results}
A well-positioned filter with preserved caval flow, no filter migration or fracture, and no pulmonary embolism during the period of protection.

\section*{Follow-up}
Periodic radiography or CT to check filter position and integrity, and plans for elective retrieval. Patients are evaluated for resuming or continuing anticoagulation and for management of the underlying thrombus.

\section*{References}
\begin{enumerate}
  \item MedlinePlus. Vena cava filter. U.S. National Library of Medicine; 2025.
  \item Current Medical Diagnosis and Treatment. McGraw-Hill; 2025.
\end{enumerate}

\clearpage
